
To continue we need to compute $\dx\left(\dfdx{y}{x}\right)$ which
seems straightforward enough since the Quotient Rule applies. This
results in
$$
\dx\left(\dfdx{y}{x}\right) =
\frac{\dx{x}\textcolor{red}{\dx{(\dx{y})}}-\dx{y}\textcolor{red}{\dx{(\dx{x})}}}{(\dx{x})^2}.
$$ 
Next simplify the notation a bit.  Let\aside{Note this very carefully,
  $\d{(\dx{y})}\neq (\dx{y})^2$, because the latter is
  $(\dx{y})\times(\dx{y})$ whereas the former is the \emph{second
    difference} of $y$.} $\textcolor{red}{\d{(\dx{y})}}=\dx^2{y}$ and
$\textcolor{red}{\d{(\dx{x})}}=\dx^2{x}$ so that,
$$
\dx\left(\dfdx{y}{x}\right) =
\frac{\dx{x}\d^2{y}-\dx{y}\dx^2{x}}{(\dx{x})^2}.  $$

But $\dx(\dx{x}) =\dx^2x$ and $\dx(\dx{y}) =\dx^2y$ present a problem
for us.  We have observed regularly beginning in
Chapter~\ref{chapt:differentials} that differentials are
questionable. If differentials are problematic, a second differential
-- the differential of a differential -- is even moreso. The time has
come to begin investigating this issue.

      %       \pagebreak[4]
\begin{Digression}[\foreign{Hic Sunt Dracones} (Here Be Dragons)]
  \label{digression:curvature-1}
  \begin{wrapfigure}[11]{r}{1.5in}
    \vskip-4mm{} \captionsetup{labelformat=empty}
    \centerline{\href{https://en.wikipedia.org/wiki/Hunt-Lenox_Globe}{\includegraphics*[height=1.5in,width=1.5in]{../Figures/LennoxGlobe2}}}
    \caption{The Lennox Globe}
    \label{fig:LennoxGlobe}
  \end{wrapfigure}
  To indicate that uncharted waters were the most dangerous medieval
  mapmakers would fill in the unexplored regions of their maps with
  illustrations of dragons and other mythological beasts. The Lennox
  Globe, in the Rare Book Division of the The New York Public Library
  (1510 CE, seen at the right), even has the inscription
  ``\foreign{Hic Sunt Dracones}'' (Here Be Dragons) along the eastern
  Asian coast.
  
  The concept of the differential of a differential, like $\dx^2x$,-
  is a mathematical dragon of sorts.  We will proceed carefully.

  The problem here is conceptual, not computational. That is, there is
  no computation we can do to see what the value of $\dx^2x$ should
  be. Rather, we need to think carefully about the meanings of our
  symbols.

  \subsubsection*{Leibniz' Conception}
  When we use Leibniz' differential notation we think of $\dx{y}$ and
  $\dx{x}$ as tiny increments in the $x$ and $y$
  coordinates. Specifically, $\dx{x}=x_2-x_1$ where $x_2$ and $x_1$
  are \emph{infinitely} close together. As we've seen the concept of a
  differential is already a difficult thing to accept.

  But if we stick with Leibniz then it must be that
  $ \dx(\dx{x}) = \dx{x_2}-\dx{x_1}, $ right? We can write this down
  and we can read these symbols easily enough. But what they seem to
  \emph{mean} is that $\dx(\dx{x})$ is the infinitely small difference
  of two infinitely small differences.

  \vskip2mm{}
  \centerline{\bf{} Down that path there are more dragons. We won't go
    there.}

  \subsubsection*{Newton's Conception}
  Instead we'll adopt Newton's viewpoint and hope that it gives us a
  meaningful interpretation of our symbols.
  
  For Newton, time was the only variable, and time always flows
  forward at a constant rate. He thought of a variable, $x$ for
  example, as moving, or flowing (\foreign{fluent}) in time and
  $\dot{x}$ was the rate of flow (\foreign{fluxion}) of the
  \foreign{fluent} $x$. In the simplest example, if $x$ is the
  position of a point then $\dot{x}$ is its velocity.

  Newton only applied his dot notation to something that was changing
  in time; something that has a \emph{rate of change with respect to
    time.} Can we apply it to time itself?

  Sure we can. Time itself is ``changing in time'', isn't it? Time
  flows at a rate of $1$ second per second, or one day per day, or one
  century per century. The units don't really matter. Loosely speaking
  (and mixing the Newtonian and Leibnizian viewpoints a bit), an
  increment of time now has the same magnitude as an increment of time
  later, has the same magnitude as an increment of time in the
  past. That is, every increment of $t$ is the same size.  If we let
  $t$ represent elapsed time then this means that the infinitesimal
  increment $\dx{t}$ is \emph{constant}, so the fluxion of $\dx{t}$ is
  zero: $\dx(\dx{t})=\dx^2{t}=\dot{\dx{t}}=0$.  However, keep in mind
  that this is not true of a quantity that is flowing at a variable
  rate. For example, if $y=t^2$ then $\eval{\dx{y}}{t}{1}=2\dx{t}$
  whereas $\eval{\dx{y}}{t}{2}=4\dx{t}$.

  This solves our problem. If $\dx{t}$ is \emph{constant} then by the
  Constant Rule $\dx^2{t}=\dx(\dx{t})=0$.

  Moreover we always have the option of adopting Newton's convention
  that the \emph{only} independent variable is time so that when we
  write $\dfdx{y}{x}$ we are thinking of the independent variable $x$
  as time (although we rarely say so out loud), and we are thinking of
  $y$ as functionally dependent on time (represented by $x$ in this
  case).  Therefore
  $$\dx^2x=\dx(\dx{x})=0 \text{ and } \dx^2y=\dx(\dx{y})\neq0.$$

  As a practical matter we will rarely be interested in second order
  differentials like $\dx^2x$ or $\dx^2y$. As we indicated in
  Section~\ref{sec:hist-intr} differentials are really just an aid to
  computation; a convenient fiction. They are a helpful aid but the
  more important and useful concept is the derivative.  Differentials
  only serve as a step along the way to finding derivatives.

  But we will be very interested in second order derivatives:
  $ \dfdx{}{x}\left(\dfdx{y}{x}\right) = \dfdxn{y}{x}{2} $ because
  these often represent real, physical quantities like velocity or
  acceleration.
  
  But now there is another difficulty that can't be ignored.  In order
  to make sense of the expression $\dx^2{x}$ we had to abandon
  Leibniz' conception altogether. It is nice that Newton's viewpoint
  gives us a way to make sense of the notation but what this tells us
  is that taking either viewpoint (Newton's or Leibniz') does not get
  us to the logical heart of the matter. A viewpoint is just a manner
  of thinking. We adopt a viewpoint in order to build intuition.

  When two viewpoints are at odds it means that we need something more
  general. In this case we need a theory that subsumes, and can
  replace, the views of both Leibniz and Newton.
  This issue must be addressed but for now we will continue to work
  intuitively.

  Eventually, when we stand at the top of the Calculus mountain we
  will be able to see the paths that Newton and Leibniz each took.  We
  will wait to begin our own trek to the summit in
  Chapter~\ref{chapt:whats-wrong-with}.
\end{Digression}

Returning to our computation in progress we have
$$
\dx\left(\dfdx{y}{x}\right) =\frac{\dx{x}\cdot
  \d^2{y}-\dx{y}\cdot\CancelToRed{=0}{\d^2{x}}}{(\dx{x})^2} =
\left(\dfdxn{y}{x}{2}\right)\dx{x}% \frac{\dx{x}\cdot\d^2{y}}{(\dx{x})^2}
$$

Returning to equation~(\ref{eq:DiffTheta}) we have
$$
\dx{\theta}=\left(\frac{1}{1+\left(\dfdx{y}{x}\right)^2}\right)\dx\hskip-2pt\left(\dfdx{y}{x}\right)=\left(\frac{1}{1+\left(\dfdx{y}{x}\right)^2}\right)\left(\dfdxn{y}{x}{2}\right)\dx{x}
$$
